% auto-generated by consolidated_report.ipynb - do not edit by hand
\begin{table}[htbp]
\centering
\caption{United States: performance over the full available window}
\label{tab:table_7_1c_us_performance_full}
\small
\resizebox{\textwidth}{!}{%
\begin{tabular}{lrrrrrrrrr}
\toprule
Portfolio & Annualised return & p-value & Sharpe ratio & p-value & Maximum drawdown & p-value & Annualised volatility & Turnover (one-way) & Effective N \\
\midrule
F-Score, equal-weight & 16.68\% & 0.605 & 0.849 & 0.569 & -38.64\% & 0.102 & 19.65\% & 62.78\% & 30.0 \\
Matched random, equal-weight & 17.02\% $\pm$ 1.16\% & --- & 0.858 $\pm$ 0.060 & --- & -41.38\% $\pm$ 2.17\% & --- & --- & --- & --- \\
F-Score, GMV (RMT-denoised) & 12.46\% & 0.113 & \textbf{0.815}$^{*}$ & 0.040 & \textbf{-29.11\%}$^{*}$ & 0.040 & 15.28\% & 64.79\% & 5.4 \\
Matched random, GMV & 10.42\% $\pm$ 1.56\% & --- & 0.642 $\pm$ 0.100 & --- & -37.55\% $\pm$ 3.70\% & --- & --- & --- & --- \\
F-Score, sector-capped GMV & \textbf{15.71\%}$^{*}$ & 0.027 & \textbf{0.957}$^{*}$ & 0.020 & -38.56\% & 0.477 & 16.41\% & 68.41\% & 8.6 \\
Matched random, sector-capped GMV & 12.43\% $\pm$ 1.60\% & --- & 0.746 $\pm$ 0.098 & --- & -38.97\% $\pm$ 3.36\% & --- & --- & --- & --- \\
\bottomrule
\end{tabular}
}
\begin{minipage}{\linewidth}\vspace{4pt}\footnotesize
Same construction as the previous table over every formation year the data support; see Table 7.0 for the two windows. Rows pair each F-Score portfolio with the matched random portfolio drawn from the same annual high-B/\allowbreak{}M universe under the same construction, so reading across a pair isolates the selection rule and reading down the pairs isolates portfolio construction. Random rows report the Monte Carlo mean $\pm$ standard deviation and carry no test of their own, being the null. Each tested statistic is followed by its one-sided empirical p-value against the matched random null: the share of Monte Carlo draws that equalled or beat the F-Score portfolio on that statistic. Maximum drawdown is a negative return, so its test treats a shallower drawdown as better. An asterisk marks significance at the 5\% level, the single level fixed in advance in §6; each statistic is starred on its own empirical p-value, not on the Sharpe ratio's. Annualised volatility, turnover and Effective N have no null attached and are therefore never starred. Returns are gross; turnover is one-way.
\end{minipage}
\end{table}
